% 04_setup.tex  --  Experimental Setup

\section{Experimental Setup}
\label{sec:setup}

\subsection{Methods Evaluated}

We evaluate 17 condensation methods spanning five groups; \Cref{tab:methods}
gives the full taxonomy. \textbf{Geometry/classical coresets} include Random
sampling, Herding~\cite{welling2009herding}, and K-center~\cite{sener2018kcenter}.
\textbf{GBDT-native subsampling} methods include GOSS~\cite{ke2017lightgbm}
and MVS~\cite{ibragimov2019mvs}, adapted to produce static condensed sets.
\textbf{Gradient- and score-based selection} methods include GraNd,
EL2N~\cite{paul2021el2n}, CRAIG-style~\cite{mirzasoleiman2020craig},
GradMatch-style~\cite{killamsetty2021gradmatch}, and gradient sampling
(loss-weighted). \textbf{Distribution-matching distillation} is represented
by DistributionMatching~\cite{zhao2023dm} adapted to tabular features.
\textbf{Gain/landscape-based methods} form our primary family: HistDistill-Greedy
(the $(1-1/e)$-optimal coverage selector), HistDistill-Refined and
HistDistill-Density (controlled landscape-fidelity diagnostic baselines), gain
sketch, gain-path sketch, and the legacy gain-path refined method. The last
three are prior gain-path methods; the paper's main claim is the diagnostic and
benchmark, not a new state-of-the-art condenser. All adaptations produce static
surrogates: GOSS/MVS select the top scored instances under the fixed GBDT state;
GraNd, EL2N, CRAIG, and GradMatch use their published scoring or matching rules;
DistributionMatching runs on processed tabular features and shares full-data
bin edges for SLD scoring.
Short names follow the original methods: GOSS is Gradient-based One-Side
Sampling, MVS is Minimal Variance Sampling, GraNd is Gradient Normed, EL2N is
Error L2 Norm, CRAIG is a gradient-matching coreset method, DM abbreviates
DistributionMatching, and HD abbreviates HistDistill.

\subsection{Datasets and Protocol}

Our experiments use public tabular datasets drawn from
OpenML~\cite{vanschoren2013openml} and the Penn Machine Learning Benchmarks
(PMLB)~\cite{olson2017pmlb}; several entries originate from UCI-style benchmark
datasets. The full inventory appears in \Cref{tab:datasets}. The repository
lists the exact OpenML/PMLB identifiers and processed split checksums. The core
study is on binary classification. For each dataset we
condense the training set to five budgets, $\{10, 25, 50, 100, 200\}$ rows, and
repeat every (dataset, method, budget) configuration over five random seeds.
Across the 26 binary datasets and the five core methods this gives $3{,}250$
retrained-model evaluations for headline performance (\Cref{tab:headline}) and
$650$ landscape-fidelity measurements, one per dataset, method, and budget.

The broader 17-method analysis uses the same 26 datasets, budgets 25/50, and
three seeds, giving 156 cells per method for \Cref{fig:universal}. The
decomposition comparison (\cref{tab:decomp}) isolates the two reference-chain
failure modes; its absolute AUROCs are not compared to the pooled headline
values. Additional studies include the $1{,}872$-cell split-noise robustness
analysis, multiclass classification, regression, hyperparameter optimisation
(HPO) transfer, and a large-scale analysis on Covertype/SUSY/HIGGS with 400-row
condensed sets; \cref{sec:limits,sec:results-scale} give details.

\subsection{Metrics}

\textbf{Downstream performance.}
Binary classification: AUROC. Multiclass: one-vs-rest macro-AUROC.
Regression: $R^2$. All scores use the model retrained on the condensed set
and evaluated on the held-out test set.

\textbf{Landscape fidelity.}
The relative $\mathrm{SLD}_\infty$ is defined in \eqref{eq:sld}.
\Cref{fig:sldlaw} uses raw (un-normalised) $\ell_\infty$ error to show
scale confounding; raw values are not comparable across datasets.
Split-regret normalises the gain cost of following the surrogate's argmax
split over the full-data optimum. Let $g^* = [\Gamma_D]_{f^*,b^*}$ be the
full-data optimal gain, and $\hat{g} = [\Gamma_D]_{\hat{f},\hat{b}}$ the gain
of the surrogate's preferred split $(\hat{f},\hat{b})=\arg\max[\Gamma_{D'}]$:
\begin{equation}
\mathrm{regret}(D,D') = \max\!\left(0,\;\frac{g^* - \hat{g}}
                                              {\max(|g^*|,\varepsilon)}\right).
\label{eq:regret}
\end{equation}

\textbf{Error decomposition.}
Three models form a reference chain: the full-data model with AUROC $A_r$
(structure reference); the condensed-set tree with leaves refit on full data,
AUROC $A_\ell$; and the condensed-set model with its own leaves, AUROC $A_d$.
\begin{equation}
\epsilon_{\mathrm{struct}} = A_r - A_\ell,\qquad
\epsilon_{\mathrm{leaf}}   = A_\ell - A_d.
\label{eq:decomp}
\end{equation}
Leaf-population mismatch $\mathrm{pop}(D')=\|p^*-p_{D'}\|_1$ measures how
representatively $D'$ fills the model leaves, where $p^*$ and $p_{D'}$ are the
empirical leaf-occupancy distributions on $D$ and $D'$.

\subsection{Learners and Reproducibility}

The reference landscape and primary downstream learner use histogram XGBoost.
Robustness checks additionally use LightGBM, CatBoost/fallback histogram
boosting, random forests, and multilayer perceptrons (MLPs). All datasets are
public and preprocessing is deterministic: categorical features follow the
repository scripts, numerical features use shared full-data bin edges, and
splits use stratified seeds where labels permit. The repository (linked in the contributions) contains seeds, hyperparameters,
package versions, processed tables, and plotting code needed to reproduce all
reported results.
